\documentclass[11pt]{article}
\usepackage[margin=1in]{geometry}
\usepackage{amsmath,amssymb,amsthm,mathtools}
\usepackage{lmodern}
\usepackage{booktabs}
\usepackage{enumitem}
\usepackage{xcolor}
\usepackage[hidelinks]{hyperref}
\hypersetup{
  pdftitle={Whole-Frame Poisson Reindexing and the Mobius-Dilation Obstruction},
  pdfauthor={Liang Wang},
  pdfsubject={Complete additive edge frames, Poisson dual packets, and a scoped obstruction}
}

\newtheorem{theorem}{Theorem}[section]
\newtheorem{proposition}[theorem]{Proposition}
\newtheorem{lemma}[theorem]{Lemma}
\newtheorem{corollary}[theorem]{Corollary}
\newtheorem{remark}[theorem]{Remark}
\newtheorem{definition}[theorem]{Definition}

\newcommand{\Fq}{\mathbb F_q}
\newcommand{\C}{\mathbb C}
\newcommand{\Z}{\mathbb Z}
\newcommand{\e}{\mathrm e}
\newcommand{\one}{\mathbf 1}
\newcommand{\Gq}{G_q}
\newcommand{\M}{\mathcal M}
\newcommand{\cE}{\mathcal E}
\newcommand{\cD}{\mathcal D}
\newcommand{\norm}[1]{\left\lVert #1\right\rVert}

\title{\textbf{Whole-Frame Poisson Reindexing\\
and the M\"obius-Dilation Obstruction}}
\author{Liang Wang\\
Huazhong University of Science and Technology\\
Wuhan 430074, P.R. China}
\date{August 2026}

\begin{document}
\emergencystretch=2em
\maketitle

\begin{abstract}
We test a proposed next transformation for the complete additive edge frame
of the standard-zero-hole reduced-residue Barban--Davenport--Halberstam
remainder.  For a fixed unit dilation $D$ modulo a prime $q$, Poisson summation
has an exact reindexing $n=qr+kD$: all edge vertices share one dual integer
lattice.  After the divisor sum is restored, however, the dual residue packet
is acted on by the multiplicative permutation $U_D:b(k)\mapsto b(kD)$.
We prove an exact vector-valued whole-frame covariance identity and diagonalize
all such permutations by multiplicative Fourier analysis.  The resulting
shared-character profile formula contains a different dual profile for every
divisor.  For the physical additive Fourier vector, a Gauss-sum calculation
returns exactly to the nonprincipal Dirichlet-character interface of the
previous scalar compiler.  Finally, an exact operator-norm calculation and a
$q=5$ resonance fixture show that arbitrary profiles can align after the
permutations.  Thus the complete frame alone cannot collapse the divisor
components into one scalar Kloosterman packet or create a power saving.  This
is a structural, scoped obstruction: the profile-aware prime-only theorem,
the strict $1/400$ gate, and the twin-prime endpoint remain open.
\end{abstract}

\section{Introduction and claim ceiling}

The current prime-dynamics route studies a fixed-gap prime-pair remainder
through a standard-zero-hole, prime-modulus, $q$-weighted reduced-residue
variance.  The previous stage constructed an exact complete-graph tight frame
on the nonzero additive frequencies and distributed the mandatory coefficient
diagonal inside the same edge cells \cite{Wang2026TPC208}.  That construction
left a specific question: can one apply the Möbius and Poisson transforms to
the complete oriented $(d,k)$ frame before paying an edge or fiber triangle,
and thereby expose one shared dual variable for a source-valid Kloosterman
estimate?

This paper answers the algebraic part of that question and records the first
failure boundary.  The answer is not a simple yes or no.  A shared dual
integer exists for each fixed divisor component.  Across divisor components,
the same dual packet is seen in different multiplicative coordinates.  The
correct whole-frame object is consequently vector-valued.  Multiplicative
Fourier analysis gives a clean normal form, but that normal form is a return to
the nonprincipal-character representation already identified in the earlier
polarized scalar compiler \cite{Wang2026V59}.

The maximum justified claim is therefore

\begin{center}
\fbox{\parbox{0.88\linewidth}{
\texttt{PROVED\_STRUCTURAL\_L1 / STOP\_SCOPED\_FRAME\_ONLY\_SAVING}.\\
The fixed-divisor Poisson reindexing, whole-frame covariance, multiplicative
profile normal form, Gauss crosswalk, and sharp alignment obstruction are
proved.  No prime-only BDH estimate, Kloosterman attachment, arithmetic
$L^2$ advance, fixed-atom credit, or twin-prime theorem is claimed.
}}
\end{center}

The distinction matters.  A finite exact frame identity is not an arithmetic
estimate, and a post-emitter Kloosterman theorem cannot be applied until the
divisor profiles, the exact $(q-2)$ diagonal subtraction, the prime shell, the
four packet signs, and the physical block reassembly have all been compiled
into its input type.

\section{The frozen additive frame}

Let $q>2$ be prime and write
\[
 \Gq=\Fq^\times,\qquad N=q-1,qquad
 P_q=I_N-\frac1N\one\one^*.
 \tag{2.1}
\]
For a finitely supported sequence $a$, a real $v$, and a scale $H>0$, define
the nonzero additive-frequency vector
\[
 Y_q[a](k;v)=\sum_{n\in\mathbb Z}a_n\e(vn/H)\e_q(-kn),
 \qquad k\in\Gq.
 \tag{2.2}
\]
If $e=\{k,l\}\subset\Gq$, the literal edge transform is
\[
 T_e[a](v)=Y_q[a](k;v)-Y_q[a](l;v).
 \tag{2.3}
\]
The complete-graph incidence vectors satisfy
\[
 \sum_{e} (e_k-e_l)(e_k-e_l)^*=NP_q.
 \tag{2.4}
\]
Consequently the edge bilinear form is
\[
 \cE_q(Y,Z):=\frac1N\sum_e (Y(k)-Y(l))
 \overline{(Z(k)-Z(l))}
 =\langle P_qY,P_qZ\rangle.
 \tag{2.5}
\]

For the TPC application, (2.5) is only the internal frame layer.  The
physical scalar retains the outer $q$, the kernel localization, the prime
shell, the exact coefficient-off-diagonal subtraction, and the signed packet
and block reassembly.  We do not replace those layers by positivity.

\section{Poisson reindexing at one divisor}

Let $D$ be an integer with $(D,q)=1$ and let $F_D$ be a Schwartz function.
Use the Fourier convention
\[
 \widehat F_D(\xi)=\int_{\mathbb R}F_D(x)\e(-x\xi)\,dx,
 \qquad \e_q(x)=\e(x/q).
 \tag{3.1}
\]
The Poisson-ready component at additive vertex $k$ is
\[
 Y_{q,D}(k)=\sum_{m\in\mathbb Z}F_D(m)\e_q(-kDm).
 \tag{3.2}
\]

\begin{theorem}[shared dual integer at fixed dilation]
\label{thm:poisson}
For every $k\in\Gq$,
\[
 \boxed{
 Y_{q,D}(k)=\sum_{r\in\mathbb Z}\widehat F_D\!\left(r+\frac{kD}{q}\right)
 =\sum_{\substack{n\in\mathbb Z\\n\equiv kD\;(q)}}\widehat F_D(n/q).}
 \tag{3.3}
\]
For fixed $D$, $(k,r)\mapsto n=qr+kD$ is a bijection from
$\Gq\times\mathbb Z$ to the integers not divisible by $q$.  If
\[
 B_{q,D}(s)=\sum_{n\equiv s\;(q)}\widehat F_D(n/q),
 \qquad s\in\Gq,
 \tag{3.4}
\]
and $(U_Db)(k)=b(kD)$, then
\[
 \boxed{Y_{q,D}=U_DB_{q,D}.}
 \tag{3.5}
\]
\end{theorem}

\begin{proof}
Poisson summation applied to $F_D(x)\e(-kDx/q)$ gives the first equality
in (3.3).  Since $D$ is invertible modulo $q$, every
$n\not\equiv0\pmod q$ has the unique representation $n=qr+kD$ with
$k\in\Gq$.  The inverse is
\[
 k\equiv nD^{-1}\pmod q,
 \qquad r=(n-kD)/q.
\]
Reindexing proves the second equality and (3.5).
\end{proof}

Thus the fixed-divisor part of the proposed route succeeds exactly: no edge
has been estimated separately, and the complete edge family uses one dual
integer lattice.  The price is the permutation $U_D$.

\section{The whole-frame vector compiler}

Let $\mathcal D$ be a finite set of unit dilations and let $c_D,d_E\in\C$.
Set
\[
 Y=\sum_{D\in\mathcal D}c_DU_DB_D,
 \qquad
 Z=\sum_{E\in\mathcal D}d_EU_EC_E.
 \tag{4.1}
\]

\begin{proposition}[whole-frame covariance]
\label{prop:covariance}
The complete edge frame gives
\[
 \boxed{
 \cE_q(Y,Z)=\sum_{D,E\in\mathcal D}c_D\overline{d_E}
 \langle P_qU_DB_D,P_qU_EC_E\rangle.}
 \tag{4.2}
\]
Moreover $U_DP_q=P_qU_D$ and the cross term is governed by the relative
dilation $D^{-1}E$:
\[
 \langle P_qU_DB_D,P_qU_EC_E\rangle
 =\langle P_qB_D,U_{D^{-1}E}P_qC_E\rangle.
 \tag{4.3}
\]
\end{proposition}

\begin{proof}
Insert (4.1) into (2.5).  A multiplicative permutation preserves the
constant vector and hence commutes with $P_q$.  Also
$U_D^*U_E=U_{D^{-1}E}$.  This proves both identities.
\end{proof}

Equation (4.2) is the exact whole-frame transform sought in the route.  In
particular, the $D\ne E$ terms are not errors.  Dropping them, or replacing
them by separate absolute values, would require a new estimate.

\section{Shared-character profile normal form}

Let $\widehat{\Gq}$ denote the multiplicative character group and define the
unitary multiplicative Fourier transform
\[
 (\M b)(\chi)=N^{-1/2}\sum_{s\in\Gq}b(s)\overline{\chi(s)}.
 \tag{5.1}
\]
Direct substitution shows
\[
 \M(U_Db)(\chi)=\chi(D)\M b(\chi).
 \tag{5.2}
\]
The principal character coordinate is the constant direction removed by
$P_q$.

\begin{theorem}[shared-character profile normal form]
\label{thm:character}
For $Y,Z$ in (4.1),
\[
 \boxed{
 \cE_q(Y,Z)=\sum_{\chi\ne\chi_0}
 \left(\sum_Dc_D\chi(D)(\M B_D)(\chi)\right)
 \overline{\left(\sum_Ed_E\chi(E)(\M C_E)(\chi)\right)}.}
 \tag{5.3}
\]
\end{theorem}

\begin{proof}
Apply the unitary transform (5.1) to (4.2), use (5.2), and delete the
principal coordinate.  Parseval gives (5.3).
\end{proof}

The exact shared coordinate is now a multiplicative character $\chi$, but
the profiles $(\M B_D)(\chi)$ still depend on $D$.  A scalar expression
$\sum_Dc_D\chi(D)$ appears only after the extra assumption that the profiles
are common.

\section{Gauss crosswalk to the previous interface}

For a nonprincipal character $\chi$, define
\[
 \tau_q(\overline\chi)=\sum_{x\in\Gq}\overline\chi(x)\e_q(x).
 \tag{6.1}
\]

\begin{proposition}[exact Gauss crosswalk]
\label{prop:gauss}
For the physical vector (2.2),
\[
 \boxed{
 (\M Y_q[a])(\chi)=
 \frac{\overline{\chi(-1)}\tau_q(\overline\chi)}{\sqrt{q-1}}
 \sum_{q\nmid n}a_n\e(vn/H)\chi(n).}
 \tag{6.2}
\]
\end{proposition}

\begin{proof}
Exchange the finite character sum and the sum over $n$.  If $q\nmid n$,
the substitution $x=-kn$ yields
\[
 \sum_{k\in\Gq}\overline\chi(k)\e_q(-kn)
 =\overline{\chi(-1)}\chi(n)\tau_q(\overline\chi).
\]
If $q\mid n$, the inner sum vanishes for nonprincipal $\chi$.  Divide by
$\sqrt{q-1}$.
\end{proof}

The right side is precisely a nonprincipal Dirichlet-character transform of
the original packet, up to the explicit Gauss factor and a harmless choice of
whether $\chi$ or $\overline\chi$ names the coordinate.  Hence the natural
whole-frame Poisson route closes back to the V59 character representation
\cite{Wang2026V59}.  It has not emitted the fixed-modulus Kloosterman arrays
accepted by the later local engine \cite{BlomerPascadi2026}.

\section{Sharp alignment obstruction}

The profile dependence in (5.3) cannot be removed by the graph identity.

\begin{theorem}[sharp vector-valued alignment]
\label{thm:alignment}
For $c=(c_D)_{D\in\mathcal D}$ define
\[
 L_c:\bigoplus_{D\in\mathcal D}\C^{\Gq}\to\C^{\Gq},
 \qquad L_c((B_D)_D)=P_q\sum_Dc_DU_DB_D.
 \tag{7.1}
\]
Then
\[
 \boxed{\norm{L_c}=\left(\sum_D|c_D|^2\right)^{1/2}.}
 \tag{7.2}
\]
Equality is attained on the centered subspace.
\end{theorem}

\begin{proof}
The adjoint is
\[
 L_c^*z=(\overline{c_D}U_D^*P_qz)_D.
\]
Therefore
\[
 L_cL_c^*z=\sum_D|c_D|^2U_DP_qU_D^*P_qz
 =\left(\sum_D|c_D|^2\right)P_qz.
\]
This proves the upper bound.  For a unit vector $z\in\one^\perp$, choose
\[
 B_D=\overline{c_D}U_D^*z/\norm{c}_2.
\]
The direct-sum input has unit norm and the output has norm $\norm{c}_2$.
\end{proof}

There is a simpler coherent fixture.  For unit weights, take
$B_D=U_D^*z$ for one nonzero centered vector $z$.  All output vectors then
align: the sum of the individual frame energies is
$|\mathcal D|\norm{z}^2$, while the whole-frame energy is
$|\mathcal D|^2\norm{z}^2$.  Thus the ratio is exactly $|\mathcal D|$.
This is an interface obstruction, not an asymptotic assertion about the
actual Möbius coefficients.

\section{Common-profile resonance and finite validation}

If one imposes the additional common-profile condition $B_D=B$, then
$M_c=\sum_Dc_DU_D$ is diagonalized by multiplicative characters:
\[
 M_c\chi=\left(\sum_Dc_D\chi(D)\right)\chi,
 \qquad
 \norm{P_qM_cP_q}=\max_{\chi\ne\chi_0}
 \left|\sum_Dc_D\chi(D)\right|.
 \tag{8.1}
\]
This special case still has exact resonance.  For $q=5$, $D=2,3$, and
$c_2=c_3=-1$, the quadratic character satisfies
$\chi(2)=\chi(3)=-1$, so its multiplier is $2$, equal to the coefficient
$\ell^1$ mass.

The project contains an independent exact certificate for
$q=3,5,7,11,13$.  It checks 1500 dual-index rows, 3016 permutation rows, all
complete-graph Laplacians, and the alignment ratios.  A separate Gaussian
sanity script checks the continuous Poisson identity in three configurations;
the maximum observed truncation error is below $10^{-12}$.  These checks are
finite QA artifacts and do not constitute asymptotic evidence.

\begin{table}[t]
\centering
\caption{Finite certificate summary.}
\label{tab:certificate}
\begin{tabular}{@{}lrl@{}}
\toprule
quantity & value & status \\
\midrule
prime moduli & $3,5,7,11,13$ & exact \\
dual reindex rows & $1500$ & exact \\
permutation-matrix rows & $3016$ & exact \\
$q=5$ coherent energy ratio & $2$ & exact \\
Gaussian Poisson cases & $3$ & numerical sanity \\
maximum Gaussian error & $4.5\times10^{-16}$ & numerical sanity \\
\bottomrule
\end{tabular}
\end{table}

The exact checks can be reproduced from the repository root with
\begin{verbatim}
cd papers/tpc-209-whole-frame-poisson-mobius-obstruction
PYTHONDONTWRITEBYTECODE=1 python -B experiments/run_certificate.py --check
PYTHONDONTWRITEBYTECODE=1 python -B experiments/independent_checker.py --check
PYTHONDONTWRITEBYTECODE=1 python -O -B experiments/independent_checker.py --check
PYTHONDONTWRITEBYTECODE=1 python -B experiments/gaussian_poisson_sanity.py
\end{verbatim}

\section{Source boundary and route decision}

The complete edge frame from TPC-208 is an exact finite-dimensional identity.
The Poisson step in Theorem~\ref{thm:poisson}, the covariance in
Proposition~\ref{prop:covariance}, the spectral form in
Theorem~\ref{thm:character}, and the Gauss crosswalk in
Proposition~\ref{prop:gauss} are proved directly here.  No external theorem
is used to promote them to an asymptotic estimate.

The current source lock contains a general-sequence BDH architecture and
post-emitter Kloosterman estimates \cite{Harper2024,BlomerPascadi2026}.  The
latter accept fixed-modulus bilinear arrays after a valid emitter and norm
ledger have been supplied.  Equations (5.3) and (6.2) are upstream of that
input: they retain divisor-dependent profiles, the exact character family,
and the physical diagonal boundary.  No locked source proves the required
prime-only, signed, block-reassembled profile bound.

The first fatal for this particular route is therefore
\[
\begin{split}
\texttt{NO\_FRAME\_ONLY\_SCALAR\_EMITTER:}
&\quad\text{Poisson gives one shared dual lattice per dilation, but}\\
&\quad\text{dilation permutations survive; character diagonalization}\\
&\quad\text{returns to V59, and arbitrary profiles align sharply.}
\end{split}
\]

This is a scoped stop, not a global nonexistence theorem.  The smallest
repaired target is a profile-aware estimate for (5.3) on the actual
Möbius/hybrid packets, preserving the $(q-2)$ subtraction, prime shell,
kernel localization, packet signs, and physical normalization.  That target
remains open and is the natural next route question.

\section{Conclusion}

The complete TPC-208 edge frame does support a legitimate whole-frame Poisson
reindexing, but only in a vector-valued sense.  The exact shared integer
$n=qr+kD$ survives all edges for one divisor.  Across divisors, the residual
permutations are not cosmetic: they are the representation-theoretic content
of the frame.  Multiplicative Fourier exposes a shared character coordinate
and simultaneously shows the exact return to the earlier V59 character
packets.  The sharp alignment theorem rules out obtaining a scalar emitter or
a power saving from frame algebra alone.

The result advances the route by replacing a vague ``look for a shared dual
variable'' instruction with a precise profile-aware interface and a precise
obstruction.  It leaves the arithmetic Gate B status unchanged:
\[
 \texttt{FULL\_GATE\_B=OPEN},\qquad
 \texttt{STRICT\_1/400=UNPAID},\qquad
 \texttt{ARITHMETIC\_ADVANCE=NO}.
\]

\bibliographystyle{plain}
\bibliography{references}

\end{document}
